\section{本章小结}
\label{sec:mrec-summary}

本章开头问：只凭有限个样本，能否知道一个未知流形在各点朝哪个方向、局部长什么样、沿它走有多远、有几个分支和几个洞。回答的前提是一个共同的采样模型。覆盖精度说明流形上有没有远离样本的地方，扰动界说明观测离开流形多远，可达半径同时度量流形在环境中的弯曲以及不同部分相互接近的程度，并保证不同分支至少相隔$2\tau$。这些量是分析的假设，不能从样本中准确读出；全域稠密也不保证每个邻域内样本在各切向上都有展开。

可达半径有两条局部推论。半径为$\tau$的球能贴着流形滚动，由此得到流形偏离切平面的二次界；对弦的中点反复投影，得到内在距离与弦长的比较。后面每一节都用到它们。

局部几何的恢复建立在切向近似上。局部PCA用邻域协方差的主方向估计切空间，它最小化到子空间的平方距离和；把弯曲与扰动写进协方差的扰动，由谱隙界得到方向误差，弯曲项随邻域半径增大、噪声项随半径减小，两者都依赖样本在切向上的展开程度。在估计的平面上用含常数项、一次项与二次项的多项式拟合法向分量，就得到连续的局部片；函数值误差不控制导数，多片的并也不自动是流形。

内在距离由邻域图上的最短路逼近。只要覆盖精度小于半径的一半、半径小于$2\tau$，图距离与内在距离之比就夹在两个因子之间，分别由$r/\tau$与$\delta/r$决定。缺采样使图距离过大甚至断开，邻接过大使它跨越相近的部分而过小，离开流形的噪声还要求半径明显大于扰动界。

拓扑的恢复走另一条路，不依赖局部拟合。样本的球并把离散点变成连续空间，Čech复形是它的神经，有限欧氏开球构成良覆盖，神经引理使二者同伦等价，同调随之可以由边界矩阵算出。在样本$\varepsilon/2$稠密且$\varepsilon<\sqrt{3/5}\,\tau$时，每一点沿直线推向最近点的路径留在球并内，球并形变收缩到流形。这恢复的是同调，不是同胚；样本离开流形或改用Rips复形都需要另行论证。

尺度未知时，持久同调让复形随尺度增长，用包含诱导的映射追踪同调类，有限过滤的持久模分解为区间，得到条形码与持久图。两张距离表逐项相差不超过$\delta$时，两套Rips过滤$\delta$交错，持久图的瓶颈距离不超过$\delta$，这是表示的稳定性。只有在恢复窗口内，横跨窗口的条才对应流形的同调；寿命长只说明跨尺度持续，不能单独证明真实孔洞的存在。

表~\ref{tab:mrec-summary}把各条线索的输入、关键条件和输出列在一起，它们各有自己的条件和结论，不能相互替代。

\begin{BookInlineTable}
  \centering
  \caption{本章各线索的条件与结论}
  \label{tab:mrec-summary}
  \begin{tabular}{lll}
    \toprule
    线索 & 关键条件 & 结论 \\
    \midrule
    切空间估计 & $\lambda_T>0$，$\beta$相对$\lambda_T$小 & 方向误差~\eqref{eq:mrec-local-pca-error} \\
    局部片拟合 & $\bm\Phi$满列秩，图表示成立 & 系数误差~\eqref{eq:mrec-coefficient-error} \\
    内在距离 & $2\delta<r<2\tau$ & 比值界~\eqref{eq:mrec-graph-lower}、\eqref{eq:mrec-graph-upper} \\
    同调恢复 & $\varepsilon/2$稠密，$\varepsilon<\sqrt{3/5}\,\tau$ & 同构~\eqref{eq:mrec-homology-recovery} \\
    持久图稳定性 & 距离表逐项相差不超过$\delta$ & $d_B\leq\delta$ \\
    \bottomrule
  \end{tabular}
\end{BookInlineTable}

本章的条件都是确定性的：给定样本与扰动界，判断结论是否成立。样本随机抽取时这些条件以多大概率满足，属于卷八的统计问题；以拟合残差或距离畸变为目标的进一步优化，则属于卷五。
